\section{From a Prescribed Force to Terminal Continuation}
\label{sec:forced-argument-overview}
\label{guide:how-proof-works}

Fix one divergence-free datum \(u_0\), one viscosity \(\nu>0\), and one
prescribed smooth force \(f\).  Let \((u,p)\) be the corresponding maximal
classical solution on \([0,T_*)\).  The proof never changes this history.
Every derivative, finite rectangle, endpoint value, and restarted solution
below belongs to the trajectory generated by the original triple
\((u_0,\nu,f)\).

The force enters the proof in two linked forms.  Its solenoidal part
\(g=\Pp f\) acts in the projected velocity equation, while its gradient part
is fixed by the normalized pressure:
\[
 \partial_tu+\Pp\nabla\!\cdot(u\otimes u)-\nu\Delta u=g,
 \qquad
 p=R_iR_j(u_i u_j)-(-\Delta)^{-1}\nabla\!\cdot f.
\]
This is why the forced argument must begin with the source.  Appending
\(f\) after constructing an unforced velocity would change both the pressure
and the history being differentiated.

The Navier--Stokes scaling acts on the complete forced equation.  If
\((u,p,f)\) is a solution, then
\[
 u_\lambda(t,x)=\lambda u(\lambda^2t,\lambda x),\qquad
 p_\lambda(t,x)=\lambda^2p(\lambda^2t,\lambda x),\qquad
 f_\lambda(t,x)=\lambda^3f(\lambda^2t,\lambda x)
\]
solves the equation with transformed datum
\(u_{0,\lambda}(x)=\lambda u_0(\lambda x)\).  For every real \(s\),
\[
 \|\Lambda^su_\lambda(t)\|_2^2
 =\lambda^{2s-1}\|\Lambda^su(\lambda^2t)\|_2^2.
\]
The invariant spatial height is therefore \(s=\tfrac12\).  This identifies
the critical row of the same forced differential structure; it does not
replace the prescribed history by its rescaling.

Write
\[
 V_n=\partial_t^nu,\qquad Q_n=\partial_t^np,
 \qquad F_n=\partial_t^nf,
\]
and let
\[
 B_n=\sum_{a=0}^n\binom na(V_a\cdot\nabla)V_{n-a}.
\]
Repeated differentiation of the original forced equation gives the complete
temporal row
\begin{align}
 V_{n+1}+B_n+\nabla Q_n&=\nu\Delta V_n+F_n,
 \label{eq:guide-forced-row}\\
 -\Delta Q_n
 &=\partial_i\partial_j\sum_{a=0}^n\binom na
      (V_a)_i(V_{n-a})_j-\nabla\!\cdot F_n.
 \label{eq:guide-forced-pressure}
\end{align}
Spatial differentiation distributes through both velocity factors and through
the prescribed source.  The fields
\[
 \partial_x^\alpha V_n,\qquad
 \partial_x^\alpha Q_n,\qquad
 \partial_x^\alpha F_n
\]
form the pressure-complete forced mixed jet.  They are coordinates of one
solution, rather than a collection of independently chosen functions.

\phantomsection\label{guide:intersection}
The mixed jet is read in two different ways.  The first is the datum-generated
compatible space--time intersection.  On every strict horizon it gives
\[
 u\in\bigcap_{r,m\ge0}H^r(0,T;H^m_\sigma(\R^3))
 \qquad\text{for every }T<T_*.
\]
The datum-generated whole-terminal intersection theorem carries these
compatible representatives to a proposed finite maximal interval
\(I=(0,T_*)\):
\[
 u\in\bigcap_{r,m\ge0}H^r(I;H^m_\sigma(\R^3)).
 \tag{A}
\]
The second reading begins independently from the critical height and the
viscous diagonal relation.  It produces finite rectangles of complete
equation rows, each carrying transport, pressure, viscosity, and the
prescribed force.  The whole-terminal intersection and the rectangle
relationships are independently derived from the same mixed jet.  Their
composition evaluates the complete rectangle family on all of \(I\).

The viscous relation is visible already in one row.  Project
\eqref{eq:guide-forced-row}, apply \(\Lambda^s\), and use
\(\Delta=-\Lambda^2\):
\[
 \Lambda^sV_{n+1}+\nu\Lambda^{s+2}V_n
 =\Lambda^s\Pp(F_n-B_n).
 \tag{B}
\]
Squaring this equality does not separate the nonlinearity from the fluid.
The cross term is the derivative of the adjacent spatial norm, so
\begin{align}
 &\|\Lambda^sV_{n+1}\|_2^2
 +\nu^2\|\Lambda^{s+2}V_n\|_2^2
 +\nu\frac{d}{dt}\|\Lambda^{s+1}V_n\|_2^2
 \notag\\
 &\hspace{4cm}
 =\|\Lambda^s\Pp(F_n-B_n)\|_2^2.
 \tag{C}
\end{align}
The complementary projection gives the pressure allocation
\[
 \nabla Q_n=(I-\Pp)(F_n-B_n),
\]
and orthogonality recombines the complete source:
\[
 \|\Lambda^s\Pp(F_n-B_n)\|_2^2
 +\|\Lambda^s\nabla Q_n\|_2^2
 =\|\Lambda^s(F_n-B_n)\|_2^2.
 \tag{D}
\]
Thus the rectangle keeps viscosity, transport, pressure, and forcing in the
same estimate.

The critical-height tower is another reading of this viscous tether.  For
\[
 E_{r,s}=\frac12\|\Lambda^sV_r\|_2^2,
\]
the complete row is
\[
 \partial_tE_{r,s}+2\nu E_{r,s+1}
 =\mathcal P_{r,s}+\mathcal W_{r,s},
 \qquad
 \mathcal W_{r,s}
 =\operatorname{Re}\langle\Lambda^sV_r,\Lambda^sF_r\rangle.
 \tag{E}
\]
At \(r=0\) and \(s=\tfrac12\), differentiating (E) repeatedly exposes
\(E_{3/2},E_{5/2},\ldots\), while the differentiated transport, pressure,
and force work remains in the same completed row.  The tower therefore
describes the geometry among mixed-jet coordinates; it does not replace the
velocity field by a scalar height.
The first two spatial heights are, explicitly,
\begin{align*}
 H'+\nu\norm{\Lambda^{3/2}u}_2^2
 &=\mathcal P_{1/2}+\mathcal W_{1/2},\\
 E_{3/2}'+\nu\norm{\Lambda^{5/2}u}_2^2
 &=\mathcal P_{3/2}+\mathcal W_{3/2},
\end{align*}
and therefore
\[
 H''=4\nu^2E_{5/2}
 -2\nu(\mathcal P_{3/2}+\mathcal W_{3/2})
 +\partial_t(\mathcal P_{1/2}+\mathcal W_{1/2}).
\]
These are the actual \(s=\tfrac12\) and \(s=\tfrac32\) rows.  Their
transport, pressure, viscosity, and force coordinates remain joined.

\phantomsection\label{guide:mixed-jet-map}
For finite \(N,M\), select the coordinates supplied by (A),
\[
 \{V_n,V_{n+1},Q_n,F_n:0\le n\le N\}
\]
at the spatial orders required by (B)--(E), and write
\[
 \mathfrak R_{N,M}(T)
 :=\sum_{n=0}^{N}\left(
 \|V_n\|_{L^2(0,T;H^{M+1})}^2
 +\|V_{n+1}\|_{L^2(0,T;H^{M-1})}^2\right).
 \tag{F}
\]
For each finite \((N,M)\), (A) first gives the finite whole-terminal sum
\[
 C_{N,M}:=\sum_{n=0}^{N}\left(
 \|V_n\|_{L^2(0,T_*;H^{M+1})}^2
 +\|V_{n+1}\|_{L^2(0,T_*;H^{M-1})}^2\right)<\infty.
\]
Restriction to \((0,T)\) and monotone exhaustion of the nonnegative norm
integrals then give
\[
 \sup_{0<T<T_*}\mathfrak R_{N,M}(T)
 =\lim_{T\uparrow T_*}\mathfrak R_{N,M}(T)
 =C_{N,M}<\infty
 \qquad(N,M<\infty).
 \tag{G}
\]
The constant may depend on the selected rectangle.  No estimate uniform in
the infinitely many indices is used.  Thus (G) is the finite projection of
the already constructed intersection, not an assumption placed on a local
solution.  The positive-sign square (C), together with (D), then evaluates
the adjacent temporal, viscous, and pressure faces on those supplied
coordinates.  Sobolev multiplication is applied within that same finite
rectangle to its binomial transport face \(B_n\); it does not replace the
nonlinear term or introduce a second premise.

Three objects meet here without being confused.  The intersection supplies
the coordinates and their compatibility.  The equations (B)--(E) supply their
viscous and critical-height relationships.  Pairing, squaring, and integrating
those relationships gives the complete-equation rectangle estimates.  The
nonlinearity remains one face of the same forced rectangle throughout.

\phantomsection\label{guide:finite-estimates}
(A) now carries the same complete equation to the terminal step.  Its
zeroth temporal projection gives, for every finite \(m\),
\[
 u\in L^2(0,T_*;H^{m+2}).
\]
Sobolev multiplication and the prescribed force then give
\begin{align*}
 \|u_t\|_{L^1(0,T_*;H^m)}
 &\le \nu\sqrt{T_*}\,\|u\|_{L^2(0,T_*;H^{m+2})}\\
 &\quad+C_m\|u\|_{L^2(0,T_*;H^{m+2})}^2
 +\|g\|_{L^1(0,T_*;H^m)}<\infty.
\end{align*}
Hence
\[
 u\in W^{1,1}(0,T_*;H^m)
 \hookrightarrow C([0,T_*];H^m)
 \qquad(m\ge0).
\]
All these traces come from the same velocity and therefore determine one
terminal state
\[
 u_*=u(T_*)\in\bigcap_{m\ge0}H^m_\sigma(\R^3).
\]

Passing \eqref{eq:guide-forced-row}--\eqref{eq:guide-forced-pressure} to
\(T_*\) determines every temporal and pressure coordinate from \(u_*\) and
the values \(F_n(T_*)\).  Local theory restarts from the condition
\(w(0)=u_*\).  In restart time \(\tau=t-T_*\), the source is
\(f(T_*+\tau)\): this is the same prescribed absolute-time force, not a new
history.  Uniqueness joins the restarted solution to the incoming one and
extends it past \(T_*\), contradicting maximality.

The rest of the paper proves the mixed-jet construction, its datum-generated
whole-terminal intersection, and the independent critical-height and rectangle
identities in their mathematical order.  Their composition carries the one
forced history through the infinite family of finite rectangle estimates, the
compatible endpoint jet, and the same-force restart.
\phantomsection\label{guide:how-proof-ends}
